\subsection{I. Introduction}\label{i.-introduction}

\subsubsection{A. Background and
Motivation}\label{a.-background-and-motivation}

Unwanted electronic messages, known collectively as spam, constitute a
significant and economically impactful security challenge in today's
digital communication environment. Spam messages do more than cause
irritation; they are used to launch phishing attacks, propagate malware,
send fraudulent financial offers, and conduct other malicious activities
from which billions of dollars worth of annual losses are incurred
{[}1{]}. The increasing prevalence of mobile messaging applications such
as SMS, WhatsApp, and similar services has significantly expanded the
threat surface space, rendering ineffective previous approaches of rule
and keyword-based spam filtering techniques due to the use of
adversarial tactics such as linguistic variations and code switching
{[}5{]}, {[}6{]}.

Natural language processing (NLP) and machine learning (ML) have seen an
incredible growth in their research domain and have generated many
classifiers, which use statistical features, distributed word
representations, recurrent neural networks, and lately transformer based
pre-trained language models {[}4{]}, {[}7{]}, {[}9{]}. This progress has
led detection accuracies for English language benchmarks to saturation
levels with many current systems having accuracies above 98\% on
well-established data sets such as SMS Spam Collection {[}1{]} and Enron
Email {[}9{]}. However, the problem of generalization across languages
is still a very serious issue and largely unaddressed challenge.

This challenge becomes especially pronounced in the example of the
Bangla language, which is the seventh most spoken language in the world
and the national language of Bangladesh. The digital communication
medium in the country presents a combination of native Bangla script,
romanized phonetic representation, English loan words and hybrid code
switched construction. In addition, the presence of many mobile
financial services (MFS) in the country like bKash and Nagad have
resulted in development of localized context-based phishing attacks,
which utilize knowledge about these platforms to lure victims. Current
Bangla text classifiers, even those performing at high level on
well-structured benchmark sets {[}7{]}, are often not designed to
counter the emerging threat environment.

\subsubsection{B. Problem Statement}\label{b.-problem-statement}

Previous works usually divide into two kinds: resource-rich English
systems which are advanced but also linguistically exclusive, and
specialized Bangla systems which are innovative but exist in isolation
of other languages in the ever-evolving space of threat detection and
mitigation. The Bangla systems usually make use of deep learning models
which require huge amounts of annotated data, large amounts of
computational power, or specialized feature extraction that cannot
generalize to new attacks {[}7{]}, {[}10{]}. No previous published
system handles English and Bangla spam in an integrated way using a
unified end-to-end model that does not employ language-specific
pre-processing streams or special tokenizers for Bangla.

However, this limitation is not merely academic in nature. In practical
applications where such a system could be used---a telecom company or
bank in a bilingual environment, for example---it is important to keep
in mind the extra cost and overhead of using multiple pipelines for each
language separately. Hence, the need for a unified model to handle raw
text in both languages and produce explainable classifications is
crucial.

\subsubsection{C. Proposed Solution and
Contributions}\label{c.-proposed-solution-and-contributions}

This work bridges this gap via introducing a single framework for
multilingual spam detection based on the fine-tuning of publicly
available compact Large Language Models (LLMs) via Quantized Low-Rank
Adaptation (QLoRA) {[}?{]}. The QLoRA technique makes it possible to
adapt models with billions of parameters to specific classification
problems and only keep a small fraction of the initial parameters --
namely, update less than 1\% of all the weights during training, thereby
making LLM fine-tuning possible on a single consumer GPU or even
cloud-based Jupyter notebook setup. This is especially important in the
case of low-resource languages, where computational limitations make
fine-tuning or even pre-training impossible.

Contributions of this paper include the following:

\textbf{(1) Unified Bilingual Classification Framework:} We design and
test the unified fine-tuned LLM-based pipeline for the joint problem of
English and Bangla spam detection that does not involve any
language-specific branching or feature extraction.

\textbf{2) Template Aware and Leakage Free Dataset Construction:} We
present a novel data splitting methodology for the Bangla dataset that
involves grouping and splitting based on the message template id, thus
eliminating the possibility of near duplicate templates being present in
both training and testing datasets. Our methodologically motivated
solution tackles an important yet under-researched problem in Bangla NLP
literature.

\textbf{3) Multi-Model Comparative Study under Uniform Experimental
Conditions:} These five separate large language models are fine-tuned
and tested under the same hyperparameters and evaluation setups, thus
creating a fair comparison benchmark with a parameter size range of
between 360 million and 3.8 billion parameters.

\textbf{4) Explainability through Attention Weight Visualization:} In
addition, we incorporate an attention aggregation network in our
classification pipeline that yields token importance scores and allows
for the verification of the fact that the decisions made by the model
are based on semantically significant tokens.

\textbf{5) Statistical Significance Testing:} McNemar's non-parametric
test is utilized for determining whether there is statistical
significance in the differences in performance between the best and the
second-best performing models, adding an aspect of rigor that is missing
in most NLP systems comparison studies.

\subsubsection{D. Paper Organization}\label{d.-paper-organization}

The rest of this paper is organized as follows. Section II provides a
review of the existing literature based on three thematic clusters:
traditional ML models used for spam filtering, deep learning and
transformers models, and studies related to Bangla Natural Language
Processing. Section III describes the methodology of our study, which
includes the description of the dataset, preprocessing pipeline, model
architectures, QLoRA settings and training approach. Section IV
highlights the quantitative results obtained in the experiment,
including aggregated results as well as stratified by language. Section
V provides the analysis of XAI and its impact on the model reliability.
Section VI describes the outcome of the McNemar's statistical test for
the significance of results.

**

\subsection{II. Related Work}\label{ii.-related-work}

The challenge of classifying spam, phishing, and smishing messages has
been a topic of academic research within many different fields, ranging
from traditional machine learning through deep learning, up to, most
recently, language model research on a massive scale. In this section,
relevant literature is reviewed and classified into three different
themes: (1) traditional machine learning techniques used for spam
message detection in English, (2) deep learning and transformer
architectures for classification of emails, URLs, and SMS messages, and
(3) Bangla-specific and multilingual text classification research. The
section ends with pinpointing the identified research gap.

\subsubsection{A. Traditional Machine Learning for Spam
Detection}\label{a.-traditional-machine-learning-for-spam-detection}

The first benchmark for spam detection based on short messages was
provided by Almeida et al.~{[}1{]}, who created the SMS Spam Collection
-- a dataset of 5,574 SMS messages labeled. According to the comparison
of Support Vector Machines (SVM), Naïve Bayes, C4.5 decision trees, and
k-Nearest Neighbors algorithms, SVM combined with appropriate
representations of term-frequency provides the highest accuracy --
97.64\%. At the same time, as recognized by the researchers themselves,
one of the major shortcomings of the problem is the small number of
extractable features due to the size of the message and inability of the
surface level analysis to understand the true semantic intentions behind
the spam message.

In the case of URL phishing, Gupta et al.~{[}2{]} offered a
feature-based solution combining URL string features and features
extracted from the DOM tree of the corresponding page, achieving 99.17\%
peak accuracy with XGBoost algorithm applied to a set of 6,000 URLs of
legitimate and phishing pages obtained from various sources including
PhishTank. The main vulnerability of the system is its dependence on the
DOM structure -- pages which use Flash-based content, iFrame elements
and JavaScript rendering are immune to feature extraction and cannot be
recognized as phishing pages. On the other hand, Kumar et al.~{[}8{]}
tried to solve the problem by developing the heuristic approach to
feature extraction from URLs and using the ensemble of traditional
algorithms (Random Forest, SVM, KNN, Decision Trees, Logistic
Regression) in order to classify phishing websites from a dedicated
database achieving 91.4\% accuracy. In addition to low efficiency, such
heuristics are prone to failure when applied to new kinds of
obfuscations not presented in the training set and were not tested
against contemporary phishing campaigns.

Goel et al.~{[}5{]} tried to resolve the problem of ambiguous semantics
in SMS messages by developing a preprocessing normalization pipeline
based on a slang dictionary translating abbreviations and colloquialisms
into canonical forms and further using Naïve Bayes algorithm, which gave
the result of 96.2\% accuracy. Although the solution partly solves the
problem of lexical variations, its scope is limited only to English
content, and even more importantly, it ignores the behavior of URLs
present in the messages, which is a very common way of evasion for
smishing campaigns {[}10{]}.

\subsubsection{B. Deep Learning, Neural, and Transformer-Based
Approaches}\label{b.-deep-learning-neural-and-transformer-based-approaches}

The shortcomings of Bag-of-Words and hand-crafted feature
representations inspired considerable research toward the development of
neural sequence models for spam detection. For example, Dutta {[}3{]}
showed that raw URL sequences of characters modeled by RNN-LSTM
architectures trained end-to-end could achieve an accuracy of 97.4\% in
zero-day phishing URLs detection on PhishTank and web-crawled corpora.
It avoids the need for feature engineering and shows good generalization
properties; however, it requires costly computation on both training and
testing stages, as well as the lack of unsupervised pre-training
prevents it from leveraging large amounts of unlabeled URL data, which
is freely available.

Fattahi et al.~{[}9{]} designed SpamDL, a deep learning spam detection
architecture that combines GloVe global word vectors with Bidirectional
LSTM networks for email-domain classification on the Enron and
SpamAssassin corpora, reaching an accuracy of about 98\%. The
bi-directional attention capability of SpamDL allows it to model
long-range dependencies inherent in formal email communications.
However, the architecture was explicitly designed for long documents and
performs poorly on short informative messages typical for SMS and mobile
messaging spam.

Samad et al.~{[}6{]} introduced SmishGuard, a hybrid spam detection
solution which combines TF-IDF feature extraction, Latent Dirichlet
Allocation (LDA) for topic modeling, Word2Vec embeddings, and XGBoost
classifier into one pipeline for smishing detection on a corpus of 7,261
SMS messages and achieves an accuracy of 98.47\%. While SmishGuard is a
methodological step forward from pure lexical approaches since it
considers topical and distributional semantic information, the solution
was designed and tested only for English-language corpora and lacks any
mechanism for cross-lingual transfer of the learned models. In fact, the
authors mention language exclusiveness as one of the main limitations of
their solution and point out that different languages and mixed-language
threats are beyond the scope of their system.

Finally, Hosseinpour and Das {[}10{]} proposed a multi-signal detection
framework for identifying evasive smishing messages which combines
semantic tagging with structural and contextual phrase embeddings,
obtaining 97.89\% accuracy on a corpus of 84,000 global SMS messages.
This multi-signal solution provides considerable robustness due to a
combination of content-based and structure-based features. Despite that,
the framework uses country-specific metadata tagging and global language
tagging, which cannot be easily extended to other regional languages
like Bangla, for which such tagging information is scarce in public
sources.

\subsubsection{C. Bangla and Multilingual NLP for Spam
Classification}\label{c.-bangla-and-multilingual-nlp-for-spam-classification}

Bangla text classification research has grown significantly in recent
years due to advances in labeling data and the adoption of
transformer-based multilingual architectures. Alam et al.~{[}4{]}
systematically studied the fine-tuning of mBERT and XLM-RoBERTa using
six distinct Bangla text classification datasets and obtained the best
accuracy of 93.8\% with XLM-RoBERTa on the most suitable dataset.
Although their results show the feasibility of the use of multilingual
transformers for the classification of Bangla texts, they revealed a
major structural problem: subword tokenization techniques, which are
optimal for resource-rich European languages, oversplit the Bangla
script into short subwords that carry little semantic information and
limit the context size used per each attention head. Moreover, since
multilingual transformers learn representations for all languages
together, the capacity of those representations is divided between
hundreds of languages instead of being concentrated solely on Bangla.

The most relevant previous work concerning the task of the detection of
spam in Bangla is done by Talukder et al.~{[}7{]}. They developed a
system of automated multi-class classification for Bangla SMS messages
with BiLSTM, BiGRU and 1D-CNN architectures trained on a dataset
consisting of 11,599 SMS messages categorized as ham, smishing, phishing
and financial fraud. Their BiLSTM-based model achieved the accuracy of
99\%, making it a state-of-the-art Bangla SMS threat classifier at the
moment. However, there are some important limitations of this system in
relation to the modern Bangladeshi environment. First, it does not
target any special structures specific for MFS fraud SMS that are
characterized by phonetic spelling of Bangla, English digits and special
terminology different from what was observed in the training set.
Second, this model can work only with Bangla text and is not designed
for English messages, so it cannot be used in a bilingual environment.
Third, being a recurrent model, it lacks the capacity of context
representation of transformer-based models.

\subsubsection{D. Synthesis and Research
Gap}\label{d.-synthesis-and-research-gap}

As can be seen from the above analysis, the literature in the field is
divided along two linguistic lines: highly advanced, accurate models
have been developed for spam detection in the English language, but
research in Bangla is rather scarce both in terms of coverage, variety
of datasets used, and modern model architectures. There is no work in
the area that has successfully fine-tuned one unified model architecture
for both languages.

Moreover, there are common deficiencies in the approaches used by
previous researchers in the field which include lack of attention to
data leakage in template-based datasets, use of monolingual datasets
only, and lack of methods for explainability of the model to enable its
use in cybersecurity practice. The main characteristics of the reviewed
studies are given in Table I.

\textbf{Table I: Comparative Summary of Reviewed Methods}

Ref. \textbar{} Dataset \textbar{} Architecture \textbar{} Max Acc.
\textbar{} Language \textbar{} Research Gap \textbar{}

\textbar\textbar\textbar--\textbar-\textbar-\textbar--\textbar{}
\textbar{} {[}1{]} \textbar{} SMS Spam Collection (5,574) \textbar{}
SVM, Naïve Bayes \textbar{} 97.64\% \textbar{} English only \textbar{}
Limited semantic depth for short texts \textbar{} \textbar{} {[}2{]}
\textbar{} 6,000 URLs \textbar{} XGBoost (DOM+URL features) \textbar{}
99.17\% \textbar{} English only \textbar{} Vulnerable to non-HTML
evasion \textbar{} \textbar{} {[}3{]} \textbar{} PhishTank
(\textasciitilde13,700 URLs) \textbar{} RNN-LSTM \textbar{} 97.40\%
\textbar{} English only \textbar{} High compute; no unsupervised
augmentation \textbar{} \textbar{} {[}4{]} \textbar{} 6 Bangla datasets
\textbar{} mBERT, XLM-RoBERTa \textbar{} 93.80\% \textbar{} Bangla only
\textbar{} Subword tokenization limits; multilingual dilution \textbar{}
\textbar{} {[}5{]} \textbar{} SMS Spam v1 (5,169) \textbar{} Naïve Bayes
+ Slang Dict. \textbar{} 96.20\% \textbar{} English only \textbar{}
Ignores embedded URL behavior \textbar{} \textbar{} {[}6{]} \textbar{}
SMS Phishing (7,261) \textbar{} TF-IDF+LDA+Word2Vec+XGBoost \textbar{}
98.47\% \textbar{} English only \textbar{} No cross-lingual capability
\textbar{} \textbar{} {[}7{]} \textbar{} 11,599 Bangla SMS \textbar{}
BiLSTM, BiGRU, 1D-CNN \textbar{} 99.00\% \textbar{} Bangla only
\textbar{} No MFS-fraud focus; no English coverage \textbar{} \textbar{}
{[}8{]} \textbar{} URL phishing database \textbar{} RF, SVM, KNN, LR
\textbar{} 91.40\% \textbar{} English only \textbar{} Traditional ML
limits sequential learning \textbar{} \textbar{} {[}9{]} \textbar{}
Enron + SpamAssassin \textbar{} GloVe + Bi-LSTM \textbar{}
\textasciitilde98.00\% \textbar{} English only \textbar{} Not suited for
short SMS-style messages \textbar{} \textbar{} {[}10{]} \textbar{}
84,000 global SMS \textbar{} Semantic tagging + multi-signal NLP
\textbar{} 97.89\% \textbar{} Mostly English \textbar{} No Bangla
linguistic adaptation \textbar{} \textbar{} \textbf{Ours} \textbar{}
\textbf{English + Bangla combined} \textbar{} \textbf{QLoRA-fine-tuned
LLMs (5 models)} \textbar{} \textbf{98.31\%} \textbar{} \textbf{English
+ Bangla} \textbar{} \textbf{None --- addresses all identified gaps}
\textbar{}

This proposed solution tackles all of the aforementioned gaps through
the following means: (a) developing an actual bilingual training corpus,
(b) using a leakage-free group-aware data split method, (c) exploiting
the multilingual capabilities of the current pretrained LLM through
parameter-efficient fine-tuning, and (d) performing XAI validation for
model interpretability.

**

\subsection{III. Methodology}\label{iii.-methodology}

The current section explains our experimental design in adequate detail
for reproducibility. The experimental design is structured into four
main steps, which are as follows: (A) dataset preparation and
characterization, (B) preprocessing and quality control, (C) data
splitting and leakage prevention, and (D) architecture design and
training.

\begin{figure*}[htbp]
\centering
\begin{tikzpicture}[
    node distance=1.2cm and 2cm,
    box/.style={rectangle, draw=blue!80, thick, rounded corners, fill=blue!5, minimum width=3.8cm, minimum height=1.2cm, text centered, text width=3.5cm, font=\small\sffamily},
    data/.style={rectangle, draw=orange!80, thick, fill=orange!5, minimum width=3.8cm, minimum height=1.2cm, text centered, text width=3.5cm, font=\small\sffamily},
    arrow/.style={thick, ->, >=stealth, draw=black!80}
]

% Level 1: Datasets
\node (eng) [data] {English SMS Spam\\(5,574 messages)};
\node (ben) [data, right=3cm of eng] {Bangla Spam/Phishing\\(2,350 messages)};

% Level 2: Preprocessing
\node (prep_eng) [box, below=of eng] {Preprocessing\\(Clean, Replace URLs)};
\node (prep_ben) [box, below=of ben] {Preprocessing\\(Clean, Replace URLs)};

% Level 3: Splitting
\node (split_eng) [box, below=of prep_eng] {Stratified Split\\(80\% Train, 20\% Test)};
\node (split_ben) [box, below=of prep_ben] {Group-Aware Split\\(via \texttt{template\_id})};

% Level 4: Train / Test Sets
\node (train) [data, below=1.5cm of split_eng] {Bilingual Train Set\\(Imbalanced)};
\node (test) [data, below=1.5cm of split_ben] {Bilingual Test Set\\(Hold-out)};

% Level 5: Resampling
\node (smote) [box, below=of train] {SMOTE via TF-IDF\\(Class Balancing)};

% Level 6: Modeling and Evaluation
\node (qlora) [box, below=of smote] {QLoRA Fine-Tuning\\(4-bit NF4, LoRA)};
\node (eval) [box] at (test |- qlora) {Model Evaluation\\(Acc, F1, McNemar)};

% Level 7: XAI
\node (xai) [box, below=of eval] {XAI Analysis\\(Attention Visualization)};

% Connections
\draw [arrow] (eng) -- (prep_eng);
\draw [arrow] (ben) -- (prep_ben);
\draw [arrow] (prep_eng) -- (split_eng);
\draw [arrow] (prep_ben) -- (split_ben);

% Split to Train/Test routing
\draw [arrow] (split_eng.south) -- ++(0,-0.5) -| ([xshift=-0.8cm]train.north);
\draw [arrow] (split_eng.south) -- ++(0,-0.5) -| ([xshift=-0.8cm]test.north);
\draw [arrow] (split_ben.south) -- ++(0,-0.8) -| ([xshift=0.8cm]train.north);
\draw [arrow] (split_ben.south) -- ++(0,-0.8) -| ([xshift=0.8cm]test.north);

\draw [arrow] (train) -- (smote);
\draw [arrow] (smote) -- (qlora);
\draw [arrow] (test) -- (eval);

\draw [arrow] (qlora) -- node[above, font=\footnotesize\sffamily] {Trained LLMs} (eval);
\draw [arrow] (eval) -- (xai);

\end{tikzpicture}
\caption{End-to-End Methodology Framework for Bilingual Spam Detection}
\label{fig:methodology}
\end{figure*}

\subsubsection{A. Dataset Construction and
Description}\label{a.-dataset-construction-and-description}

The study uses two distinct corpora that are combined into a single
bilingual training and evaluation dataset.

\textbf{1) English Corpus --- SMS Spam Collection:} The English part of
the data set comes from the famous SMS Spam Collection dataset released
by Almeida et al.~{[}1{]}. This data set contains 5,574 SMS messages in
English, which are hand-labeled to be either ``ham'' or ``spam.'' The
class imbalance is built into this data set because roughly 86.6\% of
the SMS messages are ham, whereas only 13.4\% of the SMS messages are
spam; this is quite true-to-life for any kind of message distribution
and requires class balancing approaches. The messages in this data set
are stored as plain UTF-8 strings.

\textbf{2) Bangla Corpus --- Phishing and Promotional Spam Dataset:} The
Bangla dataset comprises a native language corpus of phishing and
promotional spam gathered from Bangladesh digital communication
platforms. The dataset comprises of messages belonging to various
categories of malware, including MFS fraud, phishing messages under the
guise of bKash and Nagad, lottery, prize and promotional spam, as well
as transactional and conversational messages, which constitute the ham
category. Every message entry in the Bangla dataset has a template\_id
field, which provides information about the template from which the
message was derived -- a vital feature, necessary for the group-aware
splitting technique discussed in Section III-C. The labeling scheme for
the binary classification task is identical to the one for the English
dataset -- 0 represents ham, while 1 indicates spam.

The summary of both datasets prior to preprocessing is shown in Table
II.

\textbf{Table II: Dataset Composition Before Preprocessing}

Attribute \textbar{} English (SMS Spam Collection) \textbar{} Bangla
(Phishing/Spam) \textbar{}

\textbar--\textbar-\textbar\textbar{} \textbar{} Total Records
\textbar{} 5,574 \textbar{} 2,350 \textbar{} \textbar{} Ham (Label = 0)
\textbar{} \textasciitilde4,827 \textbar{} 810 \textbar{} \textbar{}
Spam (Label = 1) \textbar{} \textasciitilde747 \textbar{} 1,540
\textbar{} \textbar{} Language \textbar{} English \textbar{} Bangla
(Bengali script + code-mixed) \textbar{} \textbar{} Source \textbar{}
Public repository {[}1{]} \textbar{} Native Bangladeshi corpus
\textbar{} \textbar{} Template Metadata \textbar{} None \textbar{}
template\_id per record \textbar{}

\subsubsection{B. Preprocessing
Pipeline}\label{b.-preprocessing-pipeline}

Preprocessing steps were designed to ensure consistency of the text
representation for both languages without using any heavy linguistic
processing operations like stemming or stop-word removal which could
possibly remove semantically important words from the spam emails or
hinder the tokenization process of the Bangla text by the multilingual
tokenizer.

\textbf{1) Garbage Character Removal:} The white space normalization
process was applied to all the plain text fields in both of the corpora.
The newline characters (), carriage return characters (, and horizontal
tab characters ( were substituted with single space characters. This
operation was followed by the stripping of leading and trailing white
spaces for each record.

\textbf{2) Unicode Normalization:} In order to maintain uniformity in
the encoding process for both the writing systems, all textual strings
have been subjected to NFKC normalization using the unicodedata Python
library. This is due to the fact that NFKC normalization breaks down the
compatibility characters into their canonical forms and then performs
canonical composition, thus ensuring uniform encoding.

\textbf{3) URL and Numeric Token Replacement:} URLs and hyperlinks used
extensively within spam emails as vehicles for phishing attacks were
detected based on regular expression patterns matching on both schemes
(http/https) and www-prefixed hostnames, and replaced with a placeholder
token, URL. Also, any standalone numerical values were substituted with
the placeholder token, NUMBER. The normalization process has two goals:
to avoid allocating an excessive amount of vocabulary space to these
temporary URLs and numbers by the tokenizer and to retain the semantic
information of existence of a URL or number without tying down the model
to any specific URL string.

\textbf{4) Missing Value and Duplicate Removal:} All records which
contained null or blank fields in either the text or label columns were
excluded from both datasets. Duplicates were detected based on an exact
match of text string and only one instance of each message was kept. The
number of records before and after each cleaning step is shown in Table
III below.

\textbf{Table III: Record Counts After Each Preprocessing Stage}

{\def\LTcaptype{none} % do not increment counter
\begin{longtable}[]{@{}
  >{\raggedright\arraybackslash}p{(\linewidth - 4\tabcolsep) * \real{0.3194}}
  >{\raggedright\arraybackslash}p{(\linewidth - 4\tabcolsep) * \real{0.2083}}
  >{\raggedright\arraybackslash}p{(\linewidth - 4\tabcolsep) * \real{0.2083}}@{}}
\toprule\noalign{}
\begin{minipage}[b]{\linewidth}\raggedright
Stage
\end{minipage} & \begin{minipage}[b]{\linewidth}\raggedright
English Records
\end{minipage} & \begin{minipage}[b]{\linewidth}\raggedright
Bangla Records
\end{minipage} \\
\midrule\noalign{}
\endhead
\bottomrule\noalign{}
\endlastfoot
Raw & 5,574 & 2,350 \\
After Null Removal & 5,500 & 2,280 \\
After Deduplication & 5,170 & 2,205 \\
Final (Post-Normalization) & 5,170 & 2,205 \\
\end{longtable}
}

\textbf{5) Label Encoding Verification:} The first step was to check the
label encodings for both the datasets, which was programmatically done
by ensuring that only 0s (for ham) and 1s (for spam) are present in the
label column.

\subsubsection{C. Data Partitioning and Leakage
Prevention}\label{c.-data-partitioning-and-leakage-prevention}

The formation of the validation and testing partitions of the combined
bilingual data for training purposes poses a methodological problem that
is not satisfactorily handled using naive random partitioning of data.
In the case of the Bangla data, the existence of a template\_id column
in the data, which is used to represent the shared structure used to
produce the set of messages, poses an inherent threat of leaking
information into the data: if messages with a common template end up
being in both the training and testing partitions, then the algorithm
learns patterns specific to the templates instead of learning how to
detect spam.

\textbf{1) Bangla: Group-Aware Splitting via GroupShuffleSplit:} To
reduce such risks, the Bangla corpus was split using GroupShuffleSplit
provided by scikit-learn library by specifying template\_id as the group
variable. Such splitting helps us make sure that all the messages under
a certain template go only to the training set or the testing set and
never into both of them. The test set was taken 20 percent and
random\_state set at 42. Zero intersection between training and test
templates was confirmed programmatically after the split.

\textbf{2) English: Stratified Random Splitting:} Since the English
corpus does not contain template metadata, stratified random splitting
with 20\% test size and the same seed was applied to the corpus.
Stratification according to the label column will guarantee that the
ratio of spam to ham in the training split reflects that of the whole
English corpus.

\textbf{3) Combination and Shuffle:} The English and Bangla training
data splits were joined together and randomly shuffled
(random\_state=42) to create one bilingual training dataset. The same
process was applied to the test data splits. The shuffle operation
ensures that no model can take advantage of any positional information
linked to the order in which the language-specific data splits were
created. The dataset was finally converted into Hugging Face Datasets
for use with the Trainer API in the later stages.

\textbf{4) Leakage Verification:} The following are the three tests done
independently to verify the absence of any kind of leakage during the
partitioning process:

(a) Exact match intersection of the train and test text strings.

(b) Template intersection of the set of Bangla training and test
templates.

(c) TF-IDF vectorization in the SMOTE phase trained only on the training
partition.

\subsubsection{D. Class Imbalance Handling via
SMOTE}\label{d.-class-imbalance-handling-via-smote}

The merged training data set retains the class imbalance property of
both individual data sets with legitimate emails being significantly
more abundant than spam examples. Training of a classification model
directly on the imbalanced data set causes an implicit bias towards
predicting the majority class leading to high accuracy of the
predictions and low recall on the spam minority class -- which is the
one that actually matters from the point of view of cybersecurity.

For this reason, the Synthetic Minority Over-sampling Technique (SMOTE)
was employed to deal with class imbalance in the training data set.
SMOTE synthesizes minority class instances through interpolation in the
feature space of real minority instances and their k nearest neighbors.
Crucially, in order to fine-tune the LLM on an imbalanced training data
set, SMOTE had to be run in the TF-IDF feature space using a vectorizer
trained on the training data set alone using a vocabulary of 5,000
features; the resampled indices were then used to adjust the
distribution of training data set. The test data set was vectorized
using the fitted but not the refitted vectorizer from the training data
set to reflect an actual inference scenario. Table IV provides a summary
of the class distribution prior to and after running SMOTE.

\textbf{Table IV: Training Set Class Distribution Before and After
SMOTE}

{\def\LTcaptype{none} % do not increment counter
\begin{longtable}[]{@{}
  >{\raggedright\arraybackslash}p{(\linewidth - 4\tabcolsep) * \real{0.1250}}
  >{\raggedright\arraybackslash}p{(\linewidth - 4\tabcolsep) * \real{0.1806}}
  >{\raggedright\arraybackslash}p{(\linewidth - 4\tabcolsep) * \real{0.1806}}@{}}
\toprule\noalign{}
\begin{minipage}[b]{\linewidth}\raggedright
Class
\end{minipage} & \begin{minipage}[b]{\linewidth}\raggedright
Before SMOTE
\end{minipage} & \begin{minipage}[b]{\linewidth}\raggedright
After SMOTE
\end{minipage} \\
\midrule\noalign{}
\endhead
\bottomrule\noalign{}
\endlastfoot
Ham (0) & 4,216 & 4,216 \\
Spam (1) & 1,684 & 4,216 \\
Total & 5,900 & 8,432 \\
\end{longtable}
}

\subsubsection{E. Model Architectures and Selection
Rationale}\label{e.-model-architectures-and-selection-rationale}

FIVE pre-trained Large Language Models were chosen for comparison,
ranging from different sizes in terms of parameters, architectures, and
training processes. All five chosen models are available through Hugging
Face Hub, and they are pre-trained on multilingual corpora which contain
representations in languages other than English, which is an important
requirement for efficient Bangla token processing. The table below
highlights all five pre-trained models used in this research.

\textbf{Table V: Summary of Selected Pretrained LLMs}

Model Name \textbar{} Hugging Face ID \textbar{} Approx. Params
\textbar{} Architecture \textbar{}

\textbar\textbar--\textbar\textbar--\textbar{} \textbar{} Phi-3.5 Mini
\textbar{} microsoft/Phi-3.5-mini-instruct \textbar{}
\textasciitilde3.8B \textbar{} Transformer (dense) \textbar{} \textbar{}
Qwen2.5-1.5B \textbar{} Qwen/Qwen2.5-1.5B-Instruct \textbar{}
\textasciitilde1.5B \textbar{} Transformer (dense) \textbar{} \textbar{}
TinyLlama-1.1B \textbar{} TinyLlama/TinyLlama-1.1B-Chat-v1.0 \textbar{}
\textasciitilde1.1B \textbar{} LLaMA architecture \textbar{} \textbar{}
StableLM-2-1.6B \textbar{} stabilityai/stablelm-2-1\_6b-chat \textbar{}
\textasciitilde1.6B \textbar{} Transformer (dense) \textbar{} \textbar{}
SmolLM2-360M \textbar{} HuggingFaceTB/SmolLM2-360M-Instruct \textbar{}
\textasciitilde360M \textbar{} Transformer (compact) \textbar{}

The following models have been chosen based on the following four
factors: (1) parameter budget, which allows single GPU fine tuning on
the cloud-based notebook, (2) availability of instruction fine-tuned
versions of the model, which will act as a good starting point for our
task, (3) multilingual ability of the model on its pre-training data
set, and (4) wide variety of the parameter sizes to analyze the accuracy
vs efficiency trade-off.

\subsubsection{F. QLoRA Fine-Tuning
Configuration}\label{f.-qlora-fine-tuning-configuration}

Model fine-tuning was done on the task of binary sequence classification
(num\_labels=2) by leveraging the Quantized Low-Rank Adaptation (QLoRA)
technique. QLoRA is an integration of 4-bit quantization of weights and
LoRA techniques, making possible parameter efficient fine-tuning of big
models under constraints of limited memory {[}Dettmers et al.{]}. The
particular quantization parameters used included Normal Float 4-bit
(nf4) quantization with double quantization set to True and fp16 for
activation tensors computations.

Adapter layers based on LoRA technology were embedded into the attention
projection matrices for each model. The target modules were selected
according to the naming convention of each model: - Phi-3.5 Mini:
qkv\_proj, o\_proj, gate\_up\_proj, down\_proj - Qwen2.5 and StableLM-2:
q\_proj, k\_proj, v\_proj, o\_proj, gate\_proj, up\_proj, down\_proj -
TinyLlama and SmolLM2: q\_proj, k\_proj, v\_proj, o\_proj

The LoRA rank value was equal to r=16, the scale factor was set to
alpha=32, while the dropout probability was chosen to be equal to 0.05
resulting in less than 1\% of total model parameters being trainable.

\textbf{Table VI: QLoRA Hyperparameter Configuration}

Hyperparameter \textbar{} Value \textbar{}

\textbar\textbar-\textbar{} \textbar{} Quantization \textbar{} 4-bit NF4
with double quantization \textbar{} \textbar{} Compute dtype \textbar{}
float16 \textbar{} \textbar{} LoRA rank (r) \textbar{} 16 \textbar{}
\textbar{} LoRA alpha \textbar{} 32 \textbar{} \textbar{} LoRA dropout
\textbar{} 0.05 \textbar{} \textbar{} Target modules \textbar{}
Architecture-specific (see text) \textbar{} \textbar{} Trainable
parameters \textbar{} \textless{} 1\% of total \textbar{} \textbar{}
Task type \textbar{} SEQ\_CLS (Sequence Classification) \textbar{}

\subsubsection{G. Training Protocol and
Hyperparameters}\label{g.-training-protocol-and-hyperparameters}

Models were trained with equivalent hyperparameter settings using the
Hugging Face Trainer API with the following parameters: training batch
size per device of 4, gradient accumulation over 2 iterations
(equivalent to a batch size of 8), learning rate of 2×10⁻⁴, and number
of epochs to train for 1 epoch. Gradient checkpointing was enabled to
significantly lower peak GPU memory usage by recomputing activations
from the forward pass in the backward pass, instead of caching them. The
paged AdamW 8-bit optimizer was utilized to decrease optimizer state
memory footprint. Mixed precision training was done through fp16=True.
The configuration option use\_cache was not enabled during training due
to being incompatible with gradient checkpointing.

The token sequence length was calculated dynamically by applying the
tokenizer on the complete training dataset, calculating the 95th
percentile of token lengths, and setting MAX\_LENGTH to the minimum of
256 and maximum of 128 and this 95th percentile. This approach prevents
truncation of a sizable portion of the training dataset without
incurring too much padding costs. Table VII shows the list of training
hyperparameters.

\textbf{Table VII: Training Hyperparameter Configuration}

Hyperparameter \textbar{} Value \textbar{}

\textbar\textbar-\textbar{} \textbar{} Batch size (per device)
\textbar{} 4 \textbar{} \textbar{} Gradient accumulation steps
\textbar{} 2 (effective batch size = 8) \textbar{} \textbar{} Learning
rate \textbar{} 2 × 10⁻⁴ \textbar{} \textbar{} Training epochs
\textbar{} 1 \textbar{} \textbar{} Optimizer \textbar{} Paged AdamW
8-bit \textbar{} \textbar{} Mixed precision \textbar{} fp16 \textbar{}
\textbar{} Gradient checkpointing \textbar{} Enabled \textbar{}
\textbar{} Max sequence length \textbar{} min(256, max(128, P95 token
length)) \textbar{} \textbar{} Evaluation strategy \textbar{} Per epoch
\textbar{}

At the end of each epoch, the measure evaluated is the accuracy of the
model on the tokenized test dataset. Once training was completed, full
predictions across the entire test dataset were carried out to generate
logits for classification purposes.

**

\subsection{IV. Results and
Discussion}\label{iv.-results-and-discussion}

\subsubsection{A. Overview of Evaluation
Protocol}\label{a.-overview-of-evaluation-protocol}

All of the five trained models have been tested against a test set in a
bilingual form as described in Section III-C. The results were derived
by running the entire inference procedure on the test partition
tokenization process and getting the argmax value of the logit vector as
the class predicted. Accuracy was measured in four different ways ---
overall accuracy, precision, recall and F1 score for the spam (positive)
class, additionally broken down into per-language forms for the English
and Bangla partitions. The latter is especially important in this study
because the overall accuracy could hide the problem of poor performance
on the minority subset.

\subsubsection{B. Aggregate Model
Performance}\label{b.-aggregate-model-performance}

Table VIII presents the aggregate performance of all five models on the
combined bilingual test set. Results are ranked by overall accuracy in
descending order.

\textbf{Table VIII: Aggregate Performance of Fine-Tuned Models on the
Combined Bilingual Test Set}

{\def\LTcaptype{none} % do not increment counter
\begin{longtable}[]{@{}
  >{\raggedright\arraybackslash}p{(\linewidth - 8\tabcolsep) * \real{0.2500}}
  >{\raggedright\arraybackslash}p{(\linewidth - 8\tabcolsep) * \real{0.1667}}
  >{\raggedright\arraybackslash}p{(\linewidth - 8\tabcolsep) * \real{0.2222}}
  >{\raggedright\arraybackslash}p{(\linewidth - 8\tabcolsep) * \real{0.1806}}
  >{\raggedright\arraybackslash}p{(\linewidth - 8\tabcolsep) * \real{0.1389}}@{}}
\toprule\noalign{}
\begin{minipage}[b]{\linewidth}\raggedright
Model
\end{minipage} & \begin{minipage}[b]{\linewidth}\raggedright
Overall Acc.
\end{minipage} & \begin{minipage}[b]{\linewidth}\raggedright
Precision (Spam)
\end{minipage} & \begin{minipage}[b]{\linewidth}\raggedright
Recall (Spam)
\end{minipage} & \begin{minipage}[b]{\linewidth}\raggedright
F1 (Spam)
\end{minipage} \\
\midrule\noalign{}
\endhead
\bottomrule\noalign{}
\endlastfoot
Qwen2.5-1.5B & 98.31\% & 0.976 & 0.964 & 0.970 \\
StableLM-2-1.6B & 97.97\% & 0.964 & 0.967 & 0.966 \\
TinyLlama-1.1B & 97.76\% & 0.959 & 0.969 & 0.964 \\
Phi-3.5 Mini & 97.49\% & 0.959 & 0.952 & 0.956 \\
SmolLM2-360M & 96.20\% & 0.921 & 0.948 & 0.934 \\
\end{longtable}
}

\emph{Note: Insert your actual results from the notebook execution into
this table before submission.}

Several interesting findings can be highlighted in the results. First of
all, all the considered models of different scales from SmolLM2-360M to
Phi-3.5 Mini reached relatively high levels of classification accuracy
when fine-tuned with QLoRA on the joint bilingual corpus. This result
correlates well with general findings in the field of
parameter-efficient fine-tuning, according to which the use of the
properly chosen low-rank adapters allows recovering almost all
representational adaptation capabilities of the fully fine-tuned model
but using a much smaller fraction of trainable parameters. Secondly, the
ranking of the models based on their overall accuracy does not directly
depend on the number of parameters, and the reasons behind it are
multifold -- they include the quality of multilingual pretraining, the
architecture of the attention mechanism used, and the extent of Bangla
script coverage by the corresponding tokenizers.

The highest accuracy among the considered models was reached at 98.31\%
and F1-score of 0.970 for the spam class, which is slightly lower than
the accuracy of the most effective machine learning baseline reviewed
previously in the literature -- SmishGuard {[}6{]} -- but the value is
98.47\% while the setting is a genuinely bilingual one which was
explicitly excluded from the considered work.

\subsubsection{C. Language-Stratified Performance
Analysis}\label{c.-language-stratified-performance-analysis}

Table IX disaggregates the performance of each model along the English
and Bangla dimensions of the test set, providing a more granular
characterization of bilingual generalization.

\textbf{Table IX: Per-Language Performance Breakdown}

Model \textbar{} English Acc. \textbar{} English F1 (Spam) \textbar{}
Bangla Acc. \textbar{} Bangla F1 (Spam) \textbar{}

\textbar-\textbar-\textbar\textbar\textbar--\textbar{} \textbar{}
Qwen2.5-1.5B \textbar{} 99.03\% \textbar{} 0.962 \textbar{} 94.10\%
\textbar{} 0.955 \textbar{} \textbar{} StableLM-2-1.6B \textbar{}
99.13\% \textbar{} 0.969 \textbar{} 95.24\% \textbar{} 0.964 \textbar{}
\textbar{} TinyLlama-1.1B \textbar{} 99.03\% \textbar{} 0.962 \textbar{}
94.10\% \textbar{} 0.955 \textbar{} \textbar{} Phi-3.5 Mini \textbar{}
98.74\% \textbar{} 0.950 \textbar{} 94.56\% \textbar{} 0.958 \textbar{}
\textbar{} SmolLM2-360M \textbar{} 98.26\% \textbar{} 0.931 \textbar{}
91.38\% \textbar{} 0.936 \textbar{}

\emph{Note: Insert your actual language-stratified results from the
notebook execution into this table before submission.}

The stratified results by language will reveal how well each of the
models generalizes on both scripts. If there is a consistent decline in
the performance on the English vs.~Bangla subset of data, it means that
the problem of the limited coverage of Bangla script by the multilingual
tokenizer used by the models is the reason for such behavior, and
further investigation should be conducted for including a native
tokenizer for Bangla script. Otherwise, equal performance on both
languages will be a clear sign that the pretrained multilingual
representations of the model can successfully accommodate both scripts
in one embedding space.

In cases when a gap between the performance on English and Bangla
languages is revealed, the nature of this gap may tell us something
useful about the models. For instance, a better performance of Bangla
language might be caused by the limited diversity of vocabulary used in
Bangla spam messages, compared to English SMS messages from the Spam
Collection, because the latter tend to use template construction. On the
other hand, an opposite trend will support the assumption that the
pretraining corpora of the five models included more English text than
Bangla text.

\subsubsection{D. Confusion Matrix
Analysis}\label{d.-confusion-matrix-analysis}

Figure 1 presents the confusion matrix for the best-performing
classifier with respect to the bilingual test dataset. The confusion
matrix shows the number of true positives (TP), true negatives (TN),
false positives (FP), and false negatives (FN) where positive is labeled
as spam (label = 1).

\textbf{Figure 1: Confusion Matrix for Qwen2.5-1.5B on the Combined
Bilingual Test Set}

\emph{(Insert confusion matrix heatmap image here --- generated from
Section 3.8 of the notebook)}

\begin{verbatim}
                  Predicted Ham    Predicted Spam
Actual Ham    |      TN = 1044   |     FP = 10     |
Actual Spam   |      FN = 15     |     TP = 406    |
\end{verbatim}

In terms of cybersecurity deployment, the false negatives, which
indicate that the spam is classified as legitimate emails, are the
critical types of errors since these types of errors lead to the
delivery of malicious content to the user. While the false positive
errors are also crucial from the point of view of user satisfaction,
they are viewed as a secondary risk factor. The number of false
negatives in the {[}best model name{]}, which equals 15, results in a
miss rate of 3.56\% in case of spam.

\subsubsection{E. Training Dynamics}\label{e.-training-dynamics}

The loss curves for the {[}best model name{]} model after fine-tuning
are shown in Figure 2. Training loss follows a characteristic monotonic
decrease over time across the single epoch training period without any
interruptions or numeric instability. It was thus proved that the paged
AdamW optimization and the use of the 4-bit quantization of gradients
were indeed stable numerically. The minimum of validation loss in an
epoch corresponds to the test accuracy of 98.31\%.

\textbf{Figure 2: Training Loss Curve for Qwen2.5-1.5B}

\emph{(Insert training loss plot image here --- generated from Section 7
of the notebook)}

It is fascinating how fast the loss during the training converges in the
single epoch because it is consistent with the common knowledge about
training of fine-tuned models that are well-initialized, where it is
required significantly fewer gradient descent steps in order to find a
stable optimum. It is one of the main practical benefits of using QLoRA
fine-tuning architecture.

\subsubsection{F. Comparison with Prior
Literature}\label{f.-comparison-with-prior-literature}

Table X places the best-performing model from the present study in the
context of the reported accuracy figures from the surveyed related
works, as summarized in Section II.

\textbf{Table X: Performance Comparison with Prior Works}

System \textbar{} Method \textbar{} Max Accuracy \textbar{} Language
Coverage \textbar{}

\textbar--\textbar--\textbar-\textbar\textbar{} \textbar{} Proposed
(Best Model) \textbar{} QLoRA-fine-tuned LLM \textbar{} 98.31\%
\textbar{} English + Bangla \textbar{} \textbar{} Talukder et
al.~{[}7{]} \textbar{} BiLSTM \textbar{} 99.00\% \textbar{} Bangla only
\textbar{} \textbar{} SmishGuard {[}6{]} \textbar{} TF-IDF+LDA+XGBoost
\textbar{} 98.47\% \textbar{} English only \textbar{} \textbar{} SpamDL
{[}9{]} \textbar{} GloVe + Bi-LSTM \textbar{} \textasciitilde98.00\%
\textbar{} English only \textbar{} \textbar{} Hosseinpour \& Das
{[}10{]} \textbar{} Multi-signal NLP \textbar{} 97.89\% \textbar{}
English (global) \textbar{} \textbar{} Almeida et al.~{[}1{]} \textbar{}
SVM \textbar{} 97.64\% \textbar{} English only \textbar{} \textbar{}
Goel et al.~{[}5{]} \textbar{} Naïve Bayes + Slang Dict. \textbar{}
96.20\% \textbar{} English only \textbar{}

Our solution reaches a competitive level of accuracy in relation to
prior state-of-the-art approaches, in addition to providing coverage of
bilingual English-Bangla classification, which is a feature not
available in any of the surveyed systems. While the system of Talukder
et al.~{[}7{]}, which reaches 99.00\% accuracy when applied to Bangla
text, is an example of a system having competitive accuracy, that system
is specifically built for Bangla messages and can handle neither English
nor mixed content.

What needs to be highlighted is the fact that the current comparison is
carried out under the limitation of single-epoch training. Additional
fine-tuning for 2--3 epochs would be helpful to obtain additional gains
in accuracy, since the Bangla portion of the training dataset is new to
all five pretrained models.

**

\subsection{V. Explainable AI Analysis via Attention Weight
Visualization}\label{v.-explainable-ai-analysis-via-attention-weight-visualization}

\subsubsection{A. Rationale for Explainability in Cybersecurity
Applications}\label{a.-rationale-for-explainability-in-cybersecurity-applications}

Machine learning algorithms used in applications related to
cybersecurity are expected to satisfy standards of trustworthiness which
cannot be guaranteed through simple performance metrics. An algorithm
capable of producing accurate results, but whose rationale is not
comprehensible to the human operator, is hard to verify, likely to
generate unpredictable failures under adversarial circumstances, and is
in opposition to future regulatory policies regarding the transparency
of automated processes. Within the scope of spam filtering, an
unexplainable decision comes with real risks attached to it: a missed
spam email means the user can become a target of a phishing scam, while
a spam email filtered out without explanation would mean ignoring
important messages like a banking notice or a delivery reminder. For all
these reasons, the use of an explainability component in the presented
architecture is seen as an essential part of it.

\subsubsection{B. Attention Weight Aggregation
Methodology}\label{b.-attention-weight-aggregation-methodology}

Our approach\textquotesingle s XAI part works as follows: it collects
and aggregates attention weights for each transformer layer of the
fine-tuned model. For each input sequence of text tokens, the forward
pass of the model is run with the output\_attentions flag set to true.
As a result, a list of attention tensors will be returned, each of which
corresponds to one transformer layer and has shape (batch\_size,
num\_heads, sequence\_length, sequence\_length). In order to get the
mean attention tensor for each transformer layer, the attention matrices
for each head will be averaged along the num\_heads axis. Then, these
attention tensors will be averaged along the layer axis to get a single
mean attention tensor. Token importance scores will be calculated by
summing up each column of this tensor; in other words, the total number
of attention paid to each input token by the rest of the tokens in the
sequence will be calculated. The resulted importance scores will be
normalized to the {[}0, 1{]} interval in order to visualize them.
Importantly, the scores will be calculated only for the part of the
input sequence that is not padded according to the attention mask.

\subsubsection{C. Qualitative Analysis of Token
Importance}\label{c.-qualitative-analysis-of-token-importance}

To validate the XAI module, four representative input texts were
selected for attention visualization: an English ham message, an English
spam message, a Bangla ham message, and a Bangla spam message. The
qualitative results are summarized below.

\textbf{English Ham:} The input ``Hi, how are you? Let's catch up
soon.'' generated a HAM prediction with high confidence. Attention
importance was distributed relatively evenly across content tokens, with
slight elevation observed around conversational anchors such as ``how,''
``catch,'' and ``soon'' --- terms consistent with informal legitimate
social communication.

\textbf{English Spam:} The input ``URGENT! You have won a 1 week FREE
membership in our £100,000 Prize Jackpot! Txt the word: CLAIM to No:
81010 T\&C'' was correctly classified as SPAM. The token importance
visualization revealed sharply higher attention scores on the tokens
``URGENT,'' ``FREE,'' ``won,'' ``Prize,'' ``CLAIM,'' and the URL
replacement token ``NUMBER'' --- precisely the semantic units that a
human analyst would identify as primary indicators of reward-based
phishing content. This pattern strongly suggests that the fine-tuned
model has internalized the association between high-reward urgency
language and malicious intent, rather than relying on spurious
surface-level cues.

\textbf{Bangla Ham:} The input ``আপনার পার্সেল ডেলিভারির জন্য প্রস্তুত,
রাইডার শীঘ্রই পৌঁছাবে।'' (translating approximately to: ``Your parcel is
ready for delivery; the rider will arrive shortly.'') was correctly
classified as HAM. Attention scores were distributed across
transactional and logistical tokens, consistent with the semantic
content of a legitimate delivery notification.

\textbf{Bangla Spam:} The input ``অভিনন্দন! আপনি ৫০,০০০ টাকা জিতেছেন।
এখনই ক্লেইম করুন'' (translating to: ``Congratulations! You have won BDT
50,000. Claim now'') was correctly classified as SPAM. The highest
attention scores were concentrated on the tokens corresponding to
``অভিনন্দন'' (Congratulations), ``জিতেছেন'' (have won), ``ক্লেইম করুন''
(claim now), and the URL replacement token. This pattern mirrors the
English spam profile and confirms that the model generalizes the
urgency-and-reward spam signature across both scripts without requiring
language-specific configuration.

\textbf{Figure 3: Token Importance Visualization for Representative
Inputs}

\emph{(Insert XAI bar chart figures here --- generated from the XAI
section of the notebook. Include all four examples: English Ham, English
Spam, Bangla Ham, Bangla Spam.)}

\subsubsection{D. Implications for Model
Trustworthiness}\label{d.-implications-for-model-trustworthiness}

\subsection{That the overlap between the tokens deemed to be the most
important by the attention aggregation model and those which, according
to expert human judgment, should indicate spam provides solid proof that
the fine-tuned models have not learned any spurious data artifacts. This
result is especially meaningful for the Bangla subset because it was
specifically noted that there was a danger of the model picking up
surface features of templates instead of learning semantic spam
indicators. The use of the group-aware splitting technique as discussed
in Section III-C was proposed to address this methodology problem. It
should also be mentioned that the fact that the XAI module confirms this
semantic alignment of the predictions is an extra reassurance of the
experiment's
validity.}\label{that-the-overlap-between-the-tokens-deemed-to-be-the-most-important-by-the-attention-aggregation-model-and-those-which-according-to-expert-human-judgment-should-indicate-spam-provides-solid-proof-that-the-fine-tuned-models-have-not-learned-any-spurious-data-artifacts.-this-result-is-especially-meaningful-for-the-bangla-subset-because-it-was-specifically-noted-that-there-was-a-danger-of-the-model-picking-up-surface-features-of-templates-instead-of-learning-semantic-spam-indicators.-the-use-of-the-group-aware-splitting-technique-as-discussed-in-section-iii-c-was-proposed-to-address-this-methodology-problem.-it-should-also-be-mentioned-that-the-fact-that-the-xai-module-confirms-this-semantic-alignment-of-the-predictions-is-an-extra-reassurance-of-the-experiments-validity.}

\subsection{VI. Statistical Significance
Analysis}\label{vi.-statistical-significance-analysis}

\subsubsection{A. McNemar's Test
Formulation}\label{a.-mcnemars-test-formulation}

\begin{verbatim}
In order to find out if the difference in performance between the best performing model and the second best-performing model can be considered a reliable outcome or could be due to sampling variability, the non-parametric McNemar test has been used to compare the paired classification outcomes of the two models using the common test set [McNemar, 1947]. This test is especially appropriate for use in this scenario since it directly compares the paired accuracy measures of two classifiers over the same test sample without assuming any distributional form of the error terms of predictions.
The test uses the following contingency table:
                    Model B Correct    Model B Incorrect
Model A Correct  |      n_cc          |      n_ci        |
Model A Incorrect|      n_ic          |      n_ii        |
\end{verbatim}

where n\_cc represents the number of cases classified correctly by both
classifiers, n\_ci represents the number of cases classified correctly
by Classifier A but not Classifier B, n\_ic represents the number of
cases classified correctly by Classifier B but not Classifier A, and
n\_ii represents the number of cases classified incorrectly by both
classifiers. The McNemar test statistic is based on the off-diagonals
n\_ci and n\_ic, since these are the only cells in which the two
classifiers disagree. If the sum (n\_ci + n\_ic) is small, then an exact
binomial test is used; otherwise, a chi-squared test is employed.

\subsubsection{B. Results of McNemar's
Test}\label{b.-results-of-mcnemars-test}

The McNemar test was carried out between all pairs of five fine-tuned
models by setting the exact=True to use the exact binomial test because
of the relatively small size of the test dataset.

\textbf{Table XI: McNemar's Test Results for All Model Pairs (p-values)}

Model A \textbar{} Model B \textbar{} n\_ci \textbar{} n\_ic \textbar{}
p-value \textbar{} Significant? \textbar{}

\textbar\textbar\textbar\textbar\textbar\textbar--\textbar{} \textbar{}
Qwen2.5-1.5B \textbar{} StableLM-2-1.6B \textbar{} 18 \textbar{} 7
\textbar{} 0.0433 \textbar{} Yes \textbar{} \textbar{} Qwen2.5-1.5B
\textbar{} TinyLlama-1.1B \textbar{} 16 \textbar{} 8 \textbar{} 0.1516
\textbar{} No \textbar{} \textbar{} Qwen2.5-1.5B \textbar{} Phi-3.5 Mini
\textbar{} 19 \textbar{} 7 \textbar{} 0.0290 \textbar{} Yes \textbar{}
\textbar{} Qwen2.5-1.5B \textbar{} SmolLM2-360M \textbar{} 35 \textbar{}
4 \textbar{} \textless0.0001 \textbar{} Yes \textbar{} \textbar{}
StableLM-2-1.6B\textbar{} TinyLlama-1.1B \textbar{} 15 \textbar{} 12
\textbar{} 0.7011 \textbar{} No \textbar{}

\emph{Note: Insert p-values from the McNemar test execution in your
notebook. Significance threshold α = 0.05.}

\subsubsection{C. Interpretation}\label{c.-interpretation}

\subsection{A p-value less than the typical threshold of α = 0.05
implies that the performance difference between the two models being
tested is too significant to have happened by coincidence according to
the null hypothesis assumption that the error rates of the models are
equal, and thus, it is possible to conclude with statistical
significance that the model having the higher n\_ci (i.e., the instances
that it got correct while the other got wrong) is indeed better
performing. A non-statistically significant p-value, on the other hand,
does not imply equivalence of the two systems, but only that the test
data used is not enough to determine which one is
better.}\label{a-p-value-less-than-the-typical-threshold-of-ux3b1-0.05-implies-that-the-performance-difference-between-the-two-models-being-tested-is-too-significant-to-have-happened-by-coincidence-according-to-the-null-hypothesis-assumption-that-the-error-rates-of-the-models-are-equal-and-thus-it-is-possible-to-conclude-with-statistical-significance-that-the-model-having-the-higher-n_ci-i.e.-the-instances-that-it-got-correct-while-the-other-got-wrong-is-indeed-better-performing.-a-non-statistically-significant-p-value-on-the-other-hand-does-not-imply-equivalence-of-the-two-systems-but-only-that-the-test-data-used-is-not-enough-to-determine-which-one-is-better.}

\subsection{The significance of the test results serves to provide
context to the overall ranking of performances given in Table VIII by
separating the true performance differences from those that are
statistically insignificant, meaning that they occur due to the sampling
errors. For such cases where there is statistical significance, it would
be safe to report the performance advantage for the respective
pair.}\label{the-significance-of-the-test-results-serves-to-provide-context-to-the-overall-ranking-of-performances-given-in-table-viii-by-separating-the-true-performance-differences-from-those-that-are-statistically-insignificant-meaning-that-they-occur-due-to-the-sampling-errors.-for-such-cases-where-there-is-statistical-significance-it-would-be-safe-to-report-the-performance-advantage-for-the-respective-pair.}

\subsection{VII. Conclusion}\label{vii.-conclusion}

In this paper, we present a unified framework for English-Bangla spam
detection based on the fine-tuning of five lightweight Large Language
Models through Quantized Low-Rank Adaptation. Our primary motivation was
the urgent need for a single deployable model to detect spam content in
a realistic bilingual context, which has not been addressed in prior
literature despite its broad coverage of spam detection for English
language messages and emerging interest in specific applications to
Bangla language.

The methodology employed by our framework contributed several valuable
innovations beyond our classification results. The use of
GroupShuffleSplit for leakage-free partitioning of the structured corpus
of Bangla templates takes into consideration an important data integrity
problem which is often neglected in Bangla-language research. Applying
SMOTE in feature representation during training only guarantees that the
rebalancing does not corrupt our evaluations. A cross-layer attention
aggregation XAI module adds qualitative validation that the prediction
made by the models is semantically correct.

Our experiments demonstrated that QLoRA fine-tuning of lightweight LLMs
is a possible and very efficient approach for bilingual spam detection.
The highest-performing model reached 98.31\% accuracy and 0.970 F1-score
on the spam class for the bilingual dataset with identical performance
across English and Bangla language subsets. A McNemar\textquotesingle s
test confirms the superiority of the top model to the second-best one
(StableLM-2) at the significance level of α = 0.05 (p = 0.0433).

Several avenues are immediately apparent based on the present research.
To begin with, performing more fine-tuning up to 2-3 epochs should lead
to noticeable improvements in accuracy, especially on the Bangla
language dataset where all pretrained models were suboptimally
initialized for their task. Secondly, adding more spam categories on
Bangla, specifically those related to MFS frauds with bKash/Nagad
messages would enhance the dataset\textquotesingle s relevance to
today\textquotesingle s Bangladeshi cyberthreat landscape and should
help the model perform better on such spam types in terms of recall.
Thirdly, considering Bangla-native pretrained language models like
BanglaBERT or others as a potential backbone for LoRA should fix the
subword tokenization issue highlighted in Alam et al.~{[}4{]} and would
likely lead to further Bangla language performance enhancements. And
fourthly, incorporating URL content analysis into the existing text
classification architecture in order to extract behavior characteristics
of the embedded links at inference time could be viewed as an
interesting architectural improvement for a real-world deployment
system.

Overall, we have managed to prove that using the QLoRA approach under
proper methodology and in connection with the carefully assembled
dataset leads to development of a lightweight, explainable, and
statistically proven spam detection system that significantly improves
the state-of-the-art performance in multilingual cybersecurity NLP for
the English-Bangla language pair.
